% Selbständige Grammatiklektion: ausschließlich Verben mit Präpositionalobjekt
\documentclass[a4paper,11pt]{article}
\usepackage[a4paper,margin=2cm]{geometry}
\usepackage[T1]{fontenc}
\usepackage[utf8]{inputenc}
\usepackage[ngerman,english]{babel}
\usepackage{lmodern}
\usepackage{microtype}
\usepackage{xcolor}
\usepackage{parskip}
\usepackage{enumitem}
\usepackage[most]{tcolorbox}
\usepackage{titlesec}
\usepackage{array}
\usepackage{longtable}
\usepackage[table]{xcolor}
\usepackage[hidelinks]{hyperref}
\definecolor{HeaderBlueDark}{RGB}{70,110,170}
\definecolor{RowBlueLight}{RGB}{235,243,255}
\definecolor{RowBlueDark}{RGB}{220,233,250}
\definecolor{SoftGray}{RGB}{90,90,90}
\setlength{\parindent}{0pt}
\setlist[itemize]{leftmargin=1.4em,itemsep=0.35em,topsep=0.35em}
\linespread{1.06}
\renewcommand{\arraystretch}{1.38}
\setlength{\tabcolsep}{5pt}
\titleformat{\section}[block]{\Large\bfseries\color{HeaderBlueDark}}{}{0pt}{}
\titlespacing{\section}{0pt}{1.3em}{0.8em}
\newtcolorbox{exbox}{enhanced,breakable,colback=RowBlueLight,colframe=RowBlueDark,boxrule=0.4pt,arc=1.5mm,left=1.2mm,right=1.2mm,top=1mm,bottom=1mm,title={Beispiele \textnormal{(Examples)}},fonttitle=\bfseries,colbacktitle=RowBlueDark,coltitle=black}
\NewDocumentEnvironment{grammarentry}{mm+b}{%
\begin{tcolorbox}[enhanced,breakable,colback=white,colframe=HeaderBlueDark,boxrule=0.6pt,arc=1.5mm,left=1.5mm,right=1.5mm,top=1mm,bottom=1mm,colbacktitle=HeaderBlueDark,coltitle=white,fonttitle=\bfseries,title={#1\hfill\normalfont\footnotesize #2}]
#3
\end{tcolorbox}
}{}
\newcommand{\GermanText}[1]{\textbf{Deutsch: }#1\par}
\newcommand{\EnglishText}[1]{{\small\color{SoftGray}\textbf{English: }#1\par}}
\newcommand{\ExampleItem}[2]{\item \textbf{#1}\\{\small\color{SoftGray}#2}}
\title{Verben mit Präpositionalobjekt\\{\large Verbs with prepositional objects}}
\date{}
\begin{document}
\maketitle
\tableofcontents
\newpage
\section{Grundregel (Basic rule)}
\begin{grammarentry}{Präpositionalobjekt}{Verb + Präposition + Kasus}
\GermanText{Ein Präpositionalobjekt ist eine Ergänzung mit einer vom Verb bestimmten Präposition. Deshalb lernt man \textbf{Verb + Präposition + Kasus} zusammen, z.\,B. \textbf{warten auf + Akkusativ} oder \textbf{teilnehmen an + Dativ}.}
\EnglishText{A prepositional object is a complement with a preposition selected by the verb. Therefore, learn the \textbf{verb + preposition + case} together, e.g. \textbf{warten auf + accusative} or \textbf{teilnehmen an + dative}.}
\GermanText{\textbf{Akk.} = Akkusativ; \textbf{Dat.} = Dativ; \textbf{jdn.} = jemanden (Akkusativ); \textbf{jdm.} = jemandem (Dativ).}
\EnglishText{\textbf{Akk.} = accusative; \textbf{Dat.} = dative; \textbf{jdn.} = someone (accusative); \textbf{jdm.} = someone (dative).}
\end{grammarentry}

\begin{exbox}
\begin{itemize}
\ExampleItem{Maria besucht ihren Freund. (Akkusativobjekt)}{Maria visits her boyfriend. (accusative object)}
\ExampleItem{Paul hilft mir. (Dativobjekt)}{Paul helps me. (dative object)}
\ExampleItem{Die Leute warten auf den Bus. (Präpositionalobjekt)}{The people are waiting for the bus. (prepositional object)}
\end{itemize}
\end{exbox}
\section{Verben mit persönlichem Subjekt (Verbs with a person as subject)}
\rowcolors{2}{RowBlueLight}{RowBlueDark}
\begin{longtable}{>{\raggedright\arraybackslash}p{0.30\textwidth} >{\raggedright\arraybackslash}p{0.34\textwidth} >{\raggedright\arraybackslash}p{0.25\textwidth}}
\rowcolor{HeaderBlueDark}
\color{white}\textbf{Verb (verb)} & \color{white}\textbf{Rektion (government)} & \color{white}\textbf{English meaning}\\
\endfirsthead
\rowcolor{HeaderBlueDark}
\color{white}\textbf{Verb (verb)} & \color{white}\textbf{Rektion (government)} & \color{white}\textbf{English meaning}\\
\endhead
\textbf{achten / aufpassen} & auf + Akk. & pay attention / watch out for\\
\textbf{anfangen / beginnen} & mit + Dat. & start / begin with\\
\textbf{aufhören} & mit + Dat. & stop doing\\
\textbf{denken} & an + Akk. & think of\\
\textbf{nachdenken} & über + Akk. & think about\\
\textbf{fliehen / flüchten} & vor + Dat. & flee from\\
\textbf{lachen} & über + Akk. & laugh about\\
\textbf{protestieren} & gegen + Akk. & protest against\\
\textbf{rechnen} & mit + Dat. & expect / count on\\
\textbf{sorgen} & für + Akk. & provide for / take care of\\
\textbf{teilnehmen} & an + Dat. & participate in\\
\textbf{warten} & auf + Akk. & wait for\\
\end{longtable}
\begin{exbox}
\begin{itemize}
\ExampleItem{Ich denke an die Prüfung.}{I am thinking about the examination.}
\ExampleItem{Wir nehmen an dem Kurs teil.}{We participate in the course.}
\ExampleItem{Er wartet auf den Zug.}{He is waiting for the train.}
\end{itemize}
\end{exbox}
\section{Verben mit Sachen als Subjekt (Verbs with things as subjects)}
\begin{grammarentry}{Sache als Subjekt}{Hinweis}
\GermanText{In dieser Gruppe kann eine Sache das Subjekt sein. Das ist nicht dasselbe wie ein unpersönliches Subjekt.}
\EnglishText{In this group, a thing can be the subject. That is not the same as an impersonal subject.}
\end{grammarentry}
\rowcolors{2}{RowBlueLight}{RowBlueDark}
\begin{longtable}{>{\raggedright\arraybackslash}p{0.30\textwidth} >{\raggedright\arraybackslash}p{0.34\textwidth} >{\raggedright\arraybackslash}p{0.25\textwidth}}
\rowcolor{HeaderBlueDark}
\color{white}\textbf{Verb (verb)} & \color{white}\textbf{Rektion (government)} & \color{white}\textbf{English meaning}\\
\endfirsthead
\rowcolor{HeaderBlueDark}
\color{white}\textbf{Verb (verb)} & \color{white}\textbf{Rektion (government)} & \color{white}\textbf{English meaning}\\
\endhead
\textbf{gehören} & zu + Dat. & belong to\\
\textbf{passen} & zu + Dat. & match / suit\\
\textbf{riechen} & nach + Dat. & smell of\\
\textbf{schmecken} & nach + Dat. & taste of\\
\end{longtable}
\begin{exbox}
\begin{itemize}
\ExampleItem{Die Jacke passt zu meiner Hose.}{The jacket matches my trousers.}
\ExampleItem{Die Suppe schmeckt nach Knoblauch.}{The soup tastes of garlic.}
\end{itemize}
\end{exbox}
\section{Akkusativobjekt und Präpositionalobjekt (Accusative and prepositional objects)}
\rowcolors{2}{RowBlueLight}{RowBlueDark}
\begin{longtable}{>{\raggedright\arraybackslash}p{0.30\textwidth} >{\raggedright\arraybackslash}p{0.34\textwidth} >{\raggedright\arraybackslash}p{0.25\textwidth}}
\rowcolor{HeaderBlueDark}
\color{white}\textbf{Verb (verb)} & \color{white}\textbf{Rektion (government)} & \color{white}\textbf{English meaning}\\
\endfirsthead
\rowcolor{HeaderBlueDark}
\color{white}\textbf{Verb (verb)} & \color{white}\textbf{Rektion (government)} & \color{white}\textbf{English meaning}\\
\endhead
\textbf{bitten + jdn.} & um + Akk. & ask someone for\\
\textbf{einladen + jdn.} & zu + Dat. & invite someone to\\
\textbf{fragen + jdn.} & nach + Dat. & ask someone about\\
\textbf{informieren + jdn.} & über + Akk. & inform someone about\\
\end{longtable}
\begin{exbox}
\begin{itemize}
\ExampleItem{Ich bitte dich um Hilfe.}{I ask you for help.}
\ExampleItem{Sie lädt mich zu ihrer Feier ein.}{She invites me to her party.}
\ExampleItem{Er informiert mich über den Termin.}{He informs me about the appointment.}
\end{itemize}
\end{exbox}
\section{Dativobjekt und Präpositionalobjekt (Dative and prepositional objects)}
\rowcolors{2}{RowBlueLight}{RowBlueDark}
\begin{longtable}{>{\raggedright\arraybackslash}p{0.30\textwidth} >{\raggedright\arraybackslash}p{0.34\textwidth} >{\raggedright\arraybackslash}p{0.25\textwidth}}
\rowcolor{HeaderBlueDark}
\color{white}\textbf{Verb (verb)} & \color{white}\textbf{Rektion (government)} & \color{white}\textbf{English meaning}\\
\endfirsthead
\rowcolor{HeaderBlueDark}
\color{white}\textbf{Verb (verb)} & \color{white}\textbf{Rektion (government)} & \color{white}\textbf{English meaning}\\
\endhead
\textbf{antworten + jdm.} & auf + Akk. & reply to someone about\\
\textbf{danken + jdm.} & für + Akk. & thank someone for\\
\textbf{gratulieren + jdm.} & zu + Dat. & congratulate someone on\\
\textbf{helfen + jdm.} & bei + Dat. & help someone with\\
\textbf{raten + jdm.} & zu + Dat. & advise someone to\\
\end{longtable}
\begin{exbox}
\begin{itemize}
\ExampleItem{Ich danke dir für die Hilfe.}{I thank you for the help.}
\ExampleItem{Er hilft mir bei der Arbeit.}{He helps me with the work.}
\ExampleItem{Sie rät ihm zu einer Pause.}{She advises him to take a break.}
\end{itemize}
\end{exbox}
\section{Reflexive Verben mit Präpositionalobjekt (Reflexive verbs with prepositional objects)}
\begin{grammarentry}{Reflexivpronomen}{Teil der Verbkonstruktion}
\GermanText{Das Reflexivpronomen gehört zum Verbmuster, z.\,B. \textbf{ich erinnere mich}, \textbf{du erinnerst dich}, \textbf{er erinnert sich}.}
\EnglishText{The reflexive pronoun belongs to the verb pattern, e.g. \textbf{ich erinnere mich}, \textbf{du erinnerst dich}, \textbf{er erinnert sich}.}
\end{grammarentry}
\rowcolors{2}{RowBlueLight}{RowBlueDark}
\begin{longtable}{>{\raggedright\arraybackslash}p{0.30\textwidth} >{\raggedright\arraybackslash}p{0.34\textwidth} >{\raggedright\arraybackslash}p{0.25\textwidth}}
\rowcolor{HeaderBlueDark}
\color{white}\textbf{Verb (verb)} & \color{white}\textbf{Rektion (government)} & \color{white}\textbf{English meaning}\\
\endfirsthead
\rowcolor{HeaderBlueDark}
\color{white}\textbf{Verb (verb)} & \color{white}\textbf{Rektion (government)} & \color{white}\textbf{English meaning}\\
\endhead
\textbf{sich ärgern} & über + Akk. & be annoyed about\\
\textbf{sich engagieren} & für + Akk. & be committed to\\
\textbf{sich erinnern} & an + Akk. & remember\\
\textbf{sich fürchten} & vor + Dat. & be afraid of\\
\textbf{sich gewöhnen} & an + Akk. & get used to\\
\textbf{sich interessieren} & für + Akk. & be interested in\\
\textbf{sich kümmern} & um + Akk. & take care of\\
\textbf{sich sorgen} & um + Akk. & worry about\\
\textbf{sich treffen} & mit + Dat. & meet with\\
\textbf{sich verabreden} & mit + Dat. & arrange to meet\\
\textbf{sich verabschieden} & von + Dat. & say goodbye to\\
\textbf{sich verlassen} & auf + Akk. & rely on\\
\textbf{sich wenden} & an + Akk. & turn to / contact\\
\textbf{sich wundern} & über + Akk. & be surprised about\\
\end{longtable}
\begin{exbox}
\begin{itemize}
\ExampleItem{Ich erinnere mich an den Urlaub.}{I remember the holiday.}
\ExampleItem{Sie kümmert sich um ihren Vater.}{She takes care of her father.}
\ExampleItem{Wir verlassen uns auf dich.}{We rely on you.}
\end{itemize}
\end{exbox}
\section{Zwei Präpositionalobjekte (Two prepositional objects)}
\begin{grammarentry}{Zwei Ergänzungen}{nicht immer gleichzeitig nötig}
\GermanText{Die erste Präpositionalergänzung benennt häufig eine Person oder Institution und die zweite ein Thema oder einen Anlass. In einem einzelnen Satz können auch nur eine der beiden Ergänzungen auftreten.}
\EnglishText{The first prepositional complement often names a person or institution and the second a topic or reason. A sentence can also include only one of the two complements.}
\end{grammarentry}
\rowcolors{2}{RowBlueLight}{RowBlueDark}
\begin{longtable}{>{\raggedright\arraybackslash}p{0.30\textwidth} >{\raggedright\arraybackslash}p{0.34\textwidth} >{\raggedright\arraybackslash}p{0.25\textwidth}}
\rowcolor{HeaderBlueDark}
\color{white}\textbf{Verb (verb)} & \color{white}\textbf{Rektion (government)} & \color{white}\textbf{English meaning}\\
\endfirsthead
\rowcolor{HeaderBlueDark}
\color{white}\textbf{Verb (verb)} & \color{white}\textbf{Rektion (government)} & \color{white}\textbf{English meaning}\\
\endhead
\textbf{diskutieren} & mit + Dat.; über + Akk. & discuss with / about\\
\textbf{sich anmelden} & bei + Dat.; für + Akk. & register with / for\\
\textbf{sich bedanken} & bei + Dat.; für + Akk. & thank someone for\\
\textbf{sich beschweren} & bei + Dat.; über + Akk. & complain to / about\\
\textbf{sich entschuldigen} & bei + Dat.; für + Akk. & apologize to / for\\
\textbf{sich erkundigen} & bei + Dat.; nach + Dat. & inquire with / about\\
\textbf{sich streiten} & mit + Dat.; über + Akk. & argue with / about\\
\textbf{sich unterhalten} & mit + Dat.; über + Akk. & talk with / about\\
\end{longtable}
\begin{exbox}
\begin{itemize}
\ExampleItem{Ich bedanke mich bei dir für die Einladung.}{I thank you for the invitation.}
\ExampleItem{Sie erkundigt sich bei der Behörde nach den Öffnungszeiten.}{She inquires at the authority about the opening hours.}
\ExampleItem{Wir unterhalten uns mit dem Lehrer über das Buch.}{We talk with the teacher about the book.}
\end{itemize}
\end{exbox}
\section{Verschiedene Präpositionen (Alternative prepositions)}
\begin{grammarentry}{Alternativen}{Bedeutungsunterschiede}
\GermanText{Verschiedene Präpositionen drücken unterschiedliche Beziehungen aus. Die Tabelle zeigt Alternativen; nicht alle Ergänzungen müssen zugleich stehen.}
\EnglishText{Different prepositions express different relationships. The table lists alternatives; the complements do not all have to appear together.}
\end{grammarentry}
\rowcolors{2}{RowBlueLight}{RowBlueDark}
\begin{longtable}{>{\raggedright\arraybackslash}p{0.30\textwidth} >{\raggedright\arraybackslash}p{0.34\textwidth} >{\raggedright\arraybackslash}p{0.25\textwidth}}
\rowcolor{HeaderBlueDark}
\color{white}\textbf{Verb (verb)} & \color{white}\textbf{Rektion (government)} & \color{white}\textbf{English meaning}\\
\endfirsthead
\rowcolor{HeaderBlueDark}
\color{white}\textbf{Verb (verb)} & \color{white}\textbf{Rektion (government)} & \color{white}\textbf{English meaning}\\
\endhead
\textbf{arbeiten} & an + Dat. / mit + Dat. / für + Akk. & work on / with / for\\
\textbf{sich bewerben} & um + Akk. / für + Akk. & apply for\\
\textbf{sich entscheiden} & für + Akk. / gegen + Akk. & decide for / against\\
\textbf{erzählen} & von + Dat. / über + Akk. & tell about\\
\textbf{sich freuen} & auf + Akk. / über + Akk. & look forward to / be glad about\\
\textbf{sprechen} & mit + Dat. / über + Akk. / von + Dat. & speak with / about\\
\textbf{sich vorbereiten} & auf + Akk. / für + Akk. & prepare for\\
\end{longtable}
\begin{exbox}
\begin{itemize}
\ExampleItem{Ich arbeite an einem neuen Projekt.}{I am working on a new project.}
\ExampleItem{Ich arbeite mit meiner Kollegin.}{I work with my colleague.}
\ExampleItem{Ich arbeite für eine Firma.}{I work for a company.}
\ExampleItem{Ich freue mich auf das Wochenende.}{I look forward to the weekend.}
\ExampleItem{Ich freue mich über deine Nachricht.}{I am happy about your message.}
\end{itemize}
\end{exbox}
\section{Fragen und Abgrenzung (Questions and distinction)}
\begin{grammarentry}{Präposition als Teil des Verbmusters}{Objekt oder Adverbial?}
\GermanText{Bei Sachen fragt man meist mit \textbf{wo(r) + Präposition}, z.\,B. \textbf{worauf?} Bei Personen fragt man mit \textbf{Präposition + wen/wem}, z.\,B. \textbf{auf wen?}, \textbf{mit wem?}.}
\EnglishText{For things, questions usually use \textbf{wo(r) + preposition}, e.g. \textbf{worauf?}. For people, use \textbf{preposition + wen/wem}, e.g. \textbf{auf wen?}, \textbf{mit wem?}.}
\GermanText{Nicht jede Präpositionalgruppe ist ein Präpositionalobjekt: In \glqq Ich warte auf den Bus\grqq{} ist \textbf{auf den Bus} das Präpositionalobjekt. In \glqq Ich warte an der Haltestelle\grqq{} ist \textbf{an der Haltestelle} eine Ortsangabe.}
\EnglishText{Not every prepositional phrase is an object: in 'Ich warte auf den Bus,' \textbf{auf den Bus} is the prepositional object; in 'Ich warte an der Haltestelle,' \textbf{an der Haltestelle} is an adverbial of place.}
\end{grammarentry}
\begin{exbox}
\begin{itemize}
\ExampleItem{Worauf wartest du? -- Auf den Bus.}{What are you waiting for? -- For the bus.}
\ExampleItem{Auf wen wartest du? -- Auf meinen Bruder.}{Who are you waiting for? -- For my brother.}
\ExampleItem{Mit wem sprichst du? -- Mit meiner Nachbarin.}{With whom are you speaking? -- With my neighbour.}
\end{itemize}
\end{exbox}
\section{Wichtige Unterschiede (Important differences)}
\begin{grammarentry}{Bedeutung je nach Präposition}{drei häufige Paare}
\GermanText{\textbf{für jemanden sorgen} bedeutet, für dessen Wohl zu sorgen; \textbf{sich um jemanden sorgen} bedeutet, um jemanden besorgt zu sein.}
\EnglishText{\textbf{für jemanden sorgen} means to provide care for someone; \textbf{sich um jemanden sorgen} means to worry about someone.}
\GermanText{\textbf{sich auf etwas freuen} bezieht sich gewöhnlich auf Zukünftiges; \textbf{sich über etwas freuen} auf etwas Gegenwärtiges oder Vergangenes.}
\EnglishText{\textbf{sich auf etwas freuen} usually looks forward to something future; \textbf{sich über etwas freuen} expresses pleasure about something present or past.}
\GermanText{\textbf{sich um eine Stelle bewerben} ist die übliche Form. Auch \textbf{sich für eine Stelle bewerben} kommt vor; mit \textbf{sich bei einer Firma bewerben} nennt man den Arbeitgeber.}
\EnglishText{\textbf{sich um eine Stelle bewerben} is the standard form. \textbf{sich für eine Stelle bewerben} also occurs; \textbf{sich bei einer Firma bewerben} identifies the employer.}
\end{grammarentry}
\section{Zusammenfassung (Summary)}
\begin{grammarentry}{Merksatz}{für B1}
\GermanText{Lerne \textbf{Verb + Präposition + Kasus} zusammen. Falls das Verb zusätzlich ein Akkusativobjekt, Dativobjekt oder Reflexivpronomen verlangt, lerne auch diesen Bestandteil.}
\EnglishText{Learn \textbf{verb + preposition + case} together. If a verb also requires an accusative object, dative object, or reflexive pronoun, learn that part as well.}
\end{grammarentry}
\end{document}
